\subsection{Distillation and Ensemble Ablation}
\label{sec:ablation-kd}

\paragraph{Effect of knowledge distillation and target cleanliness.}
Table~\ref{tab:main} separates two questions: whether the student benefits
from a teacher at all, and whether the way we generate teacher targets matters.
In the primary configuration, \emph{distilled (TTA) targets}, the student learns
from soft labels generated by averaging each teacher's predictions over a
horizontal-flip pair of each training image (see Section~\ref{sec:kd}); it
ends at \textbf{36.6$\pm$2.1} errors.  With \emph{distilled (clean) targets},
the soft labels come from a single clean forward pass through a larger
15-teacher pool with no flip averaging, and the student reaches
\textbf{38.2$\pm$2.2} errors.  The \emph{supervised control}, trained with no
teacher signal at all, obtains \textbf{39.8$\pm$2.8} errors.

Three errors separate the best extracted Student and the supervised baseline -- moderate, and none of the three McNemar $p$ values approach $0.05$ (tables ~ \ref{tab:main}). Distillation helps, but only a little. What's more, is how unimportant the target recipe is. Students in the clean target and those in the TTA target are statistically indistinguishable (38.2 vs. 36.6, with significant seed overlap in the error distributions). A flipped average target cannot beat a clean single pass forward pass through the same teacher. A useful signal is the teacher itself. For simplicity, we keep a small 9-teacher pool with a flipped average target as the primary configuration, but the results do not depend on that choice.

\paragraph{Teacher ensemble size.} Natural follow-up: Is the plateau of 30 errors an upper limit imposed by having only 9 teachers? To find out, we ran the ensemble as \emph{direct predictor} on the test set, rather than as the source of training time soft targets, and compared 9 teachers to 15 teachers. The answer from table ~ \ref{tab:ensemble} is straightforward. More teachers do not move the ceiling.

\begin{table}[t]\centering
\caption{Teacher ensembles on DHCD test ($n{=}13{,}800$). Rows 1--2: scaling from 9 to 15 teachers does not move the 30-error plateau. Row 3: adding flip-TTA crosses $p{<}0.05$ but is methodologically invalid (Section~\ref{sec:flip-tta}).}\label{tab:ensemble}
\begin{tabular}{lcccc}\toprule
Ensemble & Errors & Acc\,\% & McNemar $p$ & $p{<}0.05$? \\\midrule
9 teachers, clean     & 30 & 99.783 & 0.0755 & no \\
15 teachers, clean    & 30 & 99.783 & 0.0755 & no \\
9 teachers, flip-TTA  & 28 & 99.797 & 0.029  & fragile \\\bottomrule
\end{tabular}\end{table}

Individual teachers do not have enough input to make the tests meaningful -- the number of clean errors per model ranges from 27 to 43. However, both majority pools reach a ceiling at \textbf{30 errors} (99.783\%), and McNemar and $p = 0.0755$ are the same. Adding 6 more teachers doesn't change anything. The remaining errors are at the non-reducible lower bound. Additional models will vote for the same incorrect answer, regardless of the number of errors it contains.

\subsection{Flip-TTA at Test Time}
\label{sec:flip-tta}

If we average each teacher's logit over the original horizontally flipped test image, the error drops to \textbf{28} ($p = 0.029$). This is the only result in our study that exceeds the 0.05 threshold. This result is not treated as a positive result. Horizontal reflection of Devanagari produces characters that are not in the script. In the training distribution, such inversions are intentionally excluded (sections ~\ref{sec:training}). Gain is also vulnerable. When applied to the \emph{individual} teacher, the flip TTA causes errors in \textbf{88} and \textbf{102} at the most affected checkpoints, increasing from 30-43 in the clean baseline. It is only by chance at this split that the ensemble masks this difference. The student training time targets (sections ~\ref{sec:kd}) also use flipped averaging, but use the \emph{training} image as the target smoothing device. This is a clear application that is not affected by the invalidity of flip TTA as a test time predictor. Therefore, we exclude the results for $p = 0.029$ and report a clean ensemble of 30 errors as an upper criterion for teacher performance.